\section{OBJETIVOS}
Os objetivos da sua pesquisa devem ser divididos em objetivo geral e específicos.

\subsection{OBJETIVO GERAL}
O objeto geral da pesquisa se refere ao tema ou assunto que está sendo investigado. É a questão principal que orienta toda a investigação e define os limites e escopo do estudo. Em outras palavras, o objeto geral da pesquisa é o foco central da investigação, o ponto de partida para a definição dos objetivos específicos e das hipóteses a serem testadas.

\subsection{OBJETIVOS ESPECÍFICOS}
Os objetivos específicos da pesquisa são metas ou etapas intermediárias que precisam ser cumpridas para alcançar o seu objetivo geral. Eles detalham e descrevem as ações que serão realizadas durante a pesquisa, e servem como guia para o pesquisador na condução do estudo. Os objetivos específicos devem ser claros, precisos e mensuráveis, e devem estar diretamente relacionados ao objetivo geral da pesquisa. É comum que cada objetivo específico tenha uma seção específica no trabalho.